\section{模型总结与拓展}
\ifteachernotes
\block{方法卡片}

长什么样：同一个顶点处有两条等长邻边夹 \(90^\circ \)，其中另两条射线夹 \(45^\circ \)；\(45^\circ \)也可能藏在两线交点处。

怎么想：先把 \(45^\circ \)搬到合适的顶点，再把分散的长度搬到同一直线上。

几步模板：确认等长与角度条件 → 转移交叉角或补出正方形 → 旋转或翻折得到全等 → 按点序写线段和差 → 用勾股或相似求值。

易错警示：矩形缺少邻边相等的条件，必须先截形或另作构造；点在延长线上时，需要重判线段加减。

常用添线：旋转适合等长邻边与半角；翻折适合半角顶点到同一直线有垂线；延长适合复制已知高，构造等腰直角三角形；矩形优先寻找局部正方形。

\block{拓展一 四个条件与定高}

在题 \(2\) 的内部构型中，四个条件任一成立，则 \(B\) 到 \(MN\) 的距离等于正方形边长 \(s\)。证明见题 \(2\)：从角平分关系得两组直角三角形全等。因此 \(S\triangle BMN=s\cdot MN/2\)。这是由斜边长求面积的快捷方法。

\block{拓展二 定周长与乘积关系}

仍用题 \(2\) 的构型，设 \(AM=m\)、\(CN=n\)、正方形边长为 \(s\)。由 \(MN=m+n\)，\(\triangle DMN\) 的周长为 \((s-m)+(s-n)+(m+n)=2s\)。

在 \(\mathrm{Rt}\triangle DMN\) 中，\((s-m)^2+(s-n)^2=(m+n)^2\)，整理得 \(mn+s(m+n)=s^2\)，即 \((s+m)(s+n)=2s^2\)。已知 \(m\) 可求 \(n=\frac{s(s-m)}{s+m}\)。该公式也可由题 \(2\) 逆向恢复半角，但使用时必须满足 \(0<m,n<s\)。

再设 \(L=MN=m+n\)。由 \(mn=s^2-sL>0\)，可得 \(L<s\)；由 \((m-n)^2=(L+2s)^2-8s^2\ge 0\)，得到 \(L\ge 2(\sqrt{2}-1)s\)。等号恰在 \(m=n=(\sqrt{2}-1)s\) 时成立。

\clearpage
\block{拓展三 底边高型的通式}

\(\triangle ABC\) 中，\(AD\perp BC\)，\(D\) 在 \(BC\) 内部，\(\angle BAC=45^\circ \)。设 \(AD=h\)、\(BD=p\)、\(CD=q\)，则 \(p,q>0\)，且 \(p,q<h\)。按题 \(5\) 的延长与旋转构造，\((p+q)^2=(h-p)^2+(h-q)^2\)，整理为 \(pq+h(p+q)=h^2\)，故 \(q=\frac{h(h-p)}{h+p}\)。反之，若 \(h,p,q>0\) 且该等式成立，同一构造可由边长恢复 \(45^\circ \)。

例如题 \(5\) 中 \(h=6\)、\(p=3\)，得 \(q=2\)；题 \(7\) 转化后 \(h=6\)、\(p=\frac{9}{4}\)，得 \(q=\frac{30}{11}\)。这里 \(D\) 必须在底边内部，不能把它直接套到垂足在外的题 \(10\) 第一种构型。

\block{拓展四 怎样判断应写和还是差}

题 \(9\) 的两问都先证明 \(MN=PN\)，再按点序表达 \(PN\)。\(P\)、\(D\)、\(N\) 时用 \(PD+DN\)；\(P\)、\(N\)、\(D\) 时用 \(PD-DN\)。解题时可以先写“目标线段＝辅助线段”，最后一步才换成题目中给出的长度。这样即使把图形转动、翻转，或把点移到延长线上，也不容易误套公式。

顶点夹半角构型中，若 \(AB=s\)、\(BE=e\)，那么 \(DF=\frac{s(s-e)}{s+e}\)，\(CF=\frac{2se}{s+e}\)，适用范围为 \(0<e<s\)。这两式把“已知 \(AB\)、\(BE\) 求 \(DF\)、\(CF\)”的手写提示补成了可直接使用的结论；其证明与拓展二的乘积关系相同。

\block{易错提醒}

\begin{itemize}
\item 矩形没有 \(AB=AD\)，不能直接套正方形的线段和。
\item 所有旋转都应注明中心、方向与角度，旋转点落在边内还是延长线上也要说明。
\item “半角”必须是两条指定射线的角，不能用交点附近任意一个 \(45^\circ \)替换。
\item 出现线段和或线段差，先判断共线点的顺序。
\item 图上看起来相等、垂直或居中，只有题目给出或证明过才能使用。
\item 求出数值后，检查长度为正、点在线段内部以及角度条件。
\end{itemize}
\else
\block{我的方法卡片}\workspace[10]
\fi
